% =====================================================================
%  Chapter 18 opener -- "Chapter at a Glance", three reading paths + the
%  ladder staircase, THE two running examples (the 16-slot array A with
%  the CSES 1648 sample in its first 8 slots, and the 10-node tree T of
%  Ch.12--17 with its Ch.16 preorder tins), the BIT-flavour cost chart,
%  the "who owns what" comparison (Ch.5 owns flattening; Ch.19/20 own
%  segment trees; Ch.1 owns merge counting; Ch.5/6 own LC 315/493/307 and
%  PBDS; Ch.17 owns PBS-with-rollback; THIS chapter owns the array-index
%  engine), recognition guide, constraint->tool table, notation.
%  Exports fig:c14-array and fig:c14-taxonomy.
%  Hand-computed values the whole chapter reuses (verified once here):
%    A=(3,2,4,5,1,1,5,3,7,0,2,4,1,6,3,2), sum 49;
%    T=(3,5,4,14,1,2,5,24,7,7,2,13,1,7,3,49);   (post-build block sums)
%    query(6)=T6+T4=16, query(7)=T7+T6+T4=21, query(13)=T13+T12+T8=38;
%    add 3 to A[3] -> T4=17,T8=27,T16=52 (query(7) then 24).
%  Every overlay below is hand-computed; every later section re-uses
%  exactly these numbers -- if a section disagrees, this opener is wrong
%  and the section is right; fix it here, never the other way.
% =====================================================================

\begin{chaptertheorybox}{Chapter at a Glance}
Chapter~5 ended every flattened tree with the sentence ``and now answer this
\emph{range} with a Fenwick tree, covered in \chref{c14}''.  Chapter~16 handed
its depth buckets to the same pointer.  This chapter is that structure: a
\term{Fenwick tree} (Binary Indexed Tree, BIT) stores, for every index \(i\),
one partial sum over the block \([i-\operatorname{lowbit}(i)+1,\,i]\) -- and
the single bit operation \(i\;\&\;(-i)\) turns both \emph{prefix queries} and
\emph{point updates} into \(\bigO{\log n}\) walks of at most \(\log n\) steps
in an ordinary array.  No nodes, no pointers, no rebalancing: the index
\emph{is} the tree.

\textbf{You will learn to}
\begin{itemize}[nosep]
  \item read responsibility intervals off the binary expansion, run the
        \code{update}/\code{query} walks, and build a BIT in \(\bigO{n}\);
        \hfill\lv{2}--\lv{3}
  \item answer point-update + range-sum end to end (CSES~1648, \(\secref{c14-core}\)); \hfill\lv{3}
  \item switch the BIT to \emph{differences} for range-add/point-read
        (CSES~1651) and pair two BITs for range-add/range-sum (\(\secref{c14-range}\)); \hfill\lv{4}
  \item use a BIT as a \emph{frequency} table: k-th smallest by binary
        lifting, inversion-style sweeps, decreasing triplets (CF~61E,
        \(\secref{c14-order}\)); \hfill\lv{4}--\lv{5}
  \item drive Ch.5's Euler array with a BIT -- subtree add/sum and the
        enter\(+1\)/leave\(-1\) root-path trick (CF~383C, \(\secref{c14-trees}\)); \hfill\lv{5}
  \item take queries \emph{offline}: sort one axis, sweep it, BIT the other
        (\(\secref{c14-offline}\)) -- dominance counting and parallel binary
        search whose predicate is a BIT instead of Ch.17's rollback DSU; and
        compose everything on the Summit CF~1181D. \hfill\lv{5}--\lv{6}
\end{itemize}

\textbf{Prerequisites.} Chapter~5 (\chref{c24}) for \(\tin/\tout\) and
``a subtree is one interval'' (\thmref{c24-range}) -- \secref{c14-trees}
consumes that reduction and never re-derives it.  Chapter~16 (\chref{c11})
for the running tree \(T\), its depths and its preorder stamps.  Chapter~17
(\chref{c28}) only for the \emph{scheduling} idea of parallel binary search,
reused by reference in \secref{c14-offline}.  Arrays, binary representation,
and \(O(\log n)\) as a concept are all that is assumed.
\end{chaptertheorybox}

% ---------------- Reading paths + ladder position ----------------
\begin{chaptertheorybox}{How to Read This Chapter (three paths)}
Level badges \(\lv1\)--\(\lv6\): \(\lv2\) is hand-feeding the lowbit walk on
a drawn array and reading intervals off the table; \(\lv3\) is a working
point-update/range-sum BIT on a real input (CSES~1648 is exactly this);
\(\lv4\) is choosing the right encoding -- differences, two BITs, or a
frequency table -- \emph{before} writing code; \(\lv5\) is composing two
chapters' machinery (Euler array \(+\) BIT; sweep \(+\) lifting) on problems
like CF~383C; \(\lv6\) is recognising that a ``static 2D query'' is really a
sorted offline sweep, or that a threshold question is PBS with a BIT
predicate -- the Summit's league.

\textbf{Ground-up path (\(\lv2\)--\(\lv4\), ``I have never written a BIT'').}
\secref{c14-core} in order: intervals, the two walks, the linear build,
CSES~1648 with every step traced on \(A\).  Then \secref{c14-range}: the
difference BIT is one change of perspective (\(\text{read}(i)=\) prefix of
differences) and the two-BIT identity is its algebraic cousin.

\textbf{Query-hammer path (\(\lv4\)--\(\lv5\), ``my TLE is a counting
problem'').}  \secref{c14-core} for the mechanics only, then
\secref{c14-order}: re-read the BIT as a table of \emph{frequencies}, where
\code{query}(k) counts elements \(\le k\).  K-th by lifting, inversion
sweeps, CF~61E.  If you only need one tool from the chapter, this section's
lifting search is it.

\textbf{Tree/offline path (\(\lv5\)--\(\lv6\), ``ranges on trees and
2D queries'').}  Ch.5's \secref{c24-flatten} (already read, by contract),
then \secref{c14-trees}: the flat array \emph{plus} this chapter's engine,
CF~383C fully traced on \(T\).  Finish with \secref{c14-offline} --
dominance sweeps and PBS-over-BIT -- and the Summit CF~1181D, which needs
all of it: offline sorting, a frequency BIT, and the lifting search
answering ``which city just crossed its quota?''.
\end{chaptertheorybox}

\begin{figure}[H]
\centering
\begin{tikzpicture}[diagram, x=9mm, y=7mm]
  \foreach \st/\lab in {1/lv1,2/lv2,3/lv3,4/lv4,5/lv5,6/lv6} {
    \fill[\lab!22, draw=\lab, line width=0.8pt] (\st-1,0) rectangle (\st,\st*0.52);
    \node[font=\scriptsize\bfseries, text=\lab] at (\st-0.5,\st*0.52-0.42) {L\st};}
  \node[annot, anchor=south west, font=\scriptsize]
    at (0.05,0.62) {18.1 lowbit + responsibility intervals (read the table)};
  \node[annot, anchor=south west, font=\scriptsize]
    at (1.05,1.14) {18.1 walks + O(n) build; solved CSES 1648};
  \node[annot, anchor=south west, font=\scriptsize]
    at (2.05,1.66) {18.2 difference BIT (CSES 1651) + two-BIT range-sum};
  \node[annot, anchor=south west, font=\scriptsize]
    at (3.05,2.18) {18.3 frequency BIT: kth-by-lifting, sweeps (CF 61E)};
  \node[annot, anchor=south west, font=\scriptsize, text=lv5]
    at (4.05,2.70) {18.4 BIT on Ch.5's Euler array: subtree/path (CF 383C)};
  \node[annot, anchor=south west, font=\scriptsize, text=lv6]
    at (5.05,3.22) {18.5 offline sweeps + PBS-over-BIT + Summit CF 1181D};
\end{tikzpicture}
\caption*{Chapter~18 on the six-step ladder: intervals and walks are
\(\lv2\)--\(\lv3\); choosing the encoding (differences, two BITs,
frequencies) is \(\lv4\); composing with Ch.5's flattening is \(\lv5\);
offline scheduling and the Summit CF~1181D are step~6.  Ground-up path
\(=\) steps 2--4; full contest path runs to~6.}
\end{figure}

% ---------------- The running examples ----------------
\subsection*{The Running Examples: One Array, One Tree}

\emph{Two} objects carry the whole chapter.  The first is a 16-slot array
whose opening half is CSES~1648's own sample, extended so that every
lowbit power up to \(16\) gets used:
\[
  A=(\,3,2,4,5,1,1,5,3,\;7,0,2,4,1,6,3,2\,),\qquad \textstyle\sum A = 49 ,
\]
with the prefix sums you will keep needing
(\(\mathit{pref}[k]=A_1+\cdots+A_k\)):
\[
  \mathit{pref}=(3,5,9,14,15,16,21,24,\;31,31,33,37,38,44,47,49).
\]
Slots \(9..16\) only ever exist to
show the \(\text{lowbit}=8\) and \(16\) intervals; every solved sample in the
chapter lives in slots \(1..8\).  The second is the tree \(T\) Chapters~12--17
have been sharing -- nodes \(1..10\), root~1, edges
\[
  1\!-\!2,\; 1\!-\!3,\; 2\!-\!4,\; 2\!-\!5,\; 3\!-\!6,\; 3\!-\!7,\; 5\!-\!8,\; 7\!-\!9,\; 9\!-\!10 .
\]
\[
  \depth=(0,1,1,2,2,2,2,3,3,4),\quad \size=(10,4,5,1,2,1,3,1,2,1),
\]
and the Ch.16 preorder (\secref{c11-paths}'s convention, \(\tout\) \emph{inclusive}):
\[
  \tin=(1,2,6,3,4,7,8,5,9,10),\qquad
  \tout=(10,5,10,3,5,7,10,5,10,10).
\]
So \(\subtree(2)=[2,5]\) holds nodes \(\{2,4,5,8\}\) and
\(\subtree(7)=[8,10]\) holds \(\{7,9,10\}\)
-- the intervals \secref{c14-trees} hands to the BIT.
\figref{c14-array} stamps \(A\) with its responsibility table and the two
walks every section reuses.

\begin{figure}[H]
\centering
\begin{tikzpicture}[diagram, scale=0.62, transform shape]
  % A row
  \foreach \i/\vv in {1/3,2/2,3/4,4/5,5/1,6/1,7/5,8/3,9/7,10/0,11/2,12/4,13/1,14/6,15/3,16/2} {
    \draw[primaryblue, line width=0.6pt, fill=primaryblue!6] (\i*0.9,0) rectangle (\i*0.9+0.9,0.9);
    \node[font=\scriptsize\ttfamily] at (\i*0.9+0.45,0.45) {\vv};
    \node[font=\tiny\ttfamily, text=darkgray] at (\i*0.9+0.45,1.15) {A[\i]};}
  % T row
  \foreach \i/\tv in {1/3,2/5,3/4,4/14,5/1,6/2,7/5,8/24,9/7,10/7,11/2,12/13,13/1,14/7,15/3,16/49} {
    \draw[difforange, line width=0.6pt, fill=difforange!8] (\i*0.9,-1.5) rectangle (\i*0.9+0.9,-0.6);
    \node[font=\scriptsize\ttfamily, text=orange!60!black] at (\i*0.9+0.45,-1.05) {\tv};
    \node[font=\tiny\ttfamily, text=darkgray] at (\i*0.9+0.45,-1.75) {T[\i]};}
  % responsibility interval bars [i-lowbit+1, i] (drawn implicitly by width below)
  \foreach \i/\lo in {1/1,2/1,3/3,4/1,5/5,6/5,7/7,8/1,9/9,10/9,11/11,12/9,13/13,14/13,15/15,16/1} {
    \fill[successgreen!30] (\lo*0.9+0.12,-2.55) rectangle (\i*0.9+0.78,-2.3);
    \node[font=\tiny, text=green!45!black] at (\i*0.9+0.45,-2.42) {[\lo,\i]};}
  \node[font=\scriptsize\ttfamily, text=primaryblue!60!black, anchor=east] at (0,0.45) {A:};
  \node[font=\scriptsize\ttfamily, text=orange!60!black, anchor=east] at (0,-1.05) {T:};
  \node[font=\scriptsize, text=green!45!black, anchor=east] at (0,-2.42) {covers:};
  % query walk 6 -> 4 -> 0
  \draw[warningred, line width=1.1pt, -latex] (6*0.9+0.45,-2.05) -- (4*0.9+0.45,-2.05);
  \draw[warningred, line width=1.1pt, -latex] (4*0.9+0.45,-2.05) -- (0*0.9+0.45,-2.05);
  \node[font=\scriptsize\ttfamily, text=warningred, anchor=north west] at (6.9,-2.1)
    {\(\text{query}(6)=T[6]+T[4]=2+14=16=\mathit{pref}[6]\)};
  % interval table
  \node[font=\scriptsize, text=darkgray, anchor=north west] at (0.3,-3.15) {%
    \(
    \begin{array}{c|cccccccccccccccc}
    i&1&2&3&4&5&6&7&8&9&10&11&12&13&14&15&16\\\hline
    \operatorname{lowbit}(i)&1&2&1&4&1&2&1&8&1&2&1&4&1&2&1&16\\
    \text{block of }T[i]&[1,1]&[1,2]&[3,3]&[1,4]&[5,5]&[5,6]&[7,7]&[1,8]&[9,9]&[9,10]&[11,11]&[9,12]&[13,13]&[13,14]&[15,15]&[1,16]
    \end{array}
    \)};
\end{tikzpicture}
\caption{The running array \(A\) (top), its BIT \(T\) (middle, \(T[i]=\) sum
of \(A\) over the green block below it), and the responsibility table
(bottom).  Build checks reused everywhere: \(T[4]=3{+}2{+}4{+}5{=}14\),
\(T[2]=3{+}2{=}5\),
\(T[8]=14{+}1{+}1{+}5{+}3{=}24\) (CSES~1648's first sample answer),
\(T[12]=7{+}0{+}2{+}4{=}13\), \(T[14]=A[13]{+}A[14]=1{+}6{=}7\), and \(T[16]=\sum A=49\).  The red
arrows show one \emph{query} walk: \(6\to4\to0\) drops the lowest set bit
each step and visits \(T[6],T[4]\) covering \([5,6]\uplus[1,4]=[1,6]\) --
disjoint blocks, no double counting.  The mirror walk (update \(A_3{+}3\))
climbs \(3\to4\to8\to16\), touching every interval that contains slot~3; both
walks are \secref{c14-core}'s opening trace.}
\label{fig:c14-array}
\end{figure}

Six angles, six hand-computed overlays -- each lives in the section named on
its row and never changes:

\begin{center}
\small
\renewcommand{\arraystretch}{1.25}
\begin{tabular}{@{}p{0.24\linewidth} p{0.50\linewidth} l@{}}
\toprule
\textbf{Angle} & \textbf{What comes out} & \textbf{Traced in} \\
\midrule
BIT over \(A\) & \(\text{query}(13){=}1{+}13{+}24{=}38{=}\mathit{pref}[13]\); add~3 to \(A_3\) cascades \(T[4],T[8],T[16]\) & \secref{c14-core} \\
difference BIT & CSES 1651's sample: \(+1\) on \([2,5]\) turns read\((4)\) from \(5\) to \(6\); read\((2)\to3\), read\((7)=5\) (outside) & \secref{c14-range} \\
two BITs & \(a{=}(3,2,4,5,1,1,5,3)\): \(\text{rangeSum}(1,4)=14\); after \(+1\) on \([2,5]\), \(\text{rangeSum}(1,4)=17\) & \secref{c14-range} \\
frequency lifting & hotels-as-counts \((3,2,4,1,5,5,2,6)\), prefix \((3,5,9,10,15,\dots)\): first prefix \(\ge6\) is slot~3; \(\ge11\) is slot~5 & \secref{c14-order} \\
sweep on \((10,8,3,1)\) & smaller-after counts \((3,2,1,0)\); inversions \(=6\); decreasing triplets \(=4\) & \secref{c14-order} \\
BIT on \(T\)'s Euler array & subtree(2) \(=[2,5]\) add \(+3\), whole tree \(+2\): node~4 (inside) reads \(5\), node~6 (outside) reads \(2\) & \secref{c14-trees} \\
\bottomrule
\end{tabular}
\end{center}

% ---------------- Which flavour, cost chart ----------------
\subsection*{Which BIT Flavour? (the cost chart)}

All five are the \emph{same} array-of-block-sums; they differ in what you
\emph{store} in slot \(i\), and that choice is forced by the query shape, not
by taste.

\begin{center}
\small
\renewcommand{\arraystretch}{1.25}
\begin{tabular}{@{}p{0.19\linewidth} p{0.15\linewidth} p{0.15\linewidth} p{0.39\linewidth}@{}}
\toprule
\textbf{Flavour} & \textbf{Update} & \textbf{Query} & \textbf{When it is the right tool} \\
\midrule
plain (18.1) & point \(\bigO{\log}\) & prefix/range sum \(\bigO{\log}\) & values change one index at a time; sums asked \\
difference (18.2) & \textbf{range} add \(\bigO{\log}\) & point read \(\bigO{\log}\) & whole intervals shift; you only ever read single positions \\
two BITs (18.2) & range add \(\bigO{\log}\) & \textbf{range} sum \(\bigO{\log}\) & both intervals move \emph{and} interval sums are asked \\
frequency (18.3) & insert/delete at rank & count \(\le k\), k-th by lifting & elements are \emph{piles of equal items}, not positions; order questions \\
on a tree (18.4) & any flavour & over \([\tin(v),\tout(v)]\) & Ch.5 already flattened the tree; now the ranges get an engine \\
\bottomrule
\end{tabular}
\end{center}

\begin{insightbox}[One array, and the trick is always ``which block owns index \(i\)'']
Every flavour in the table is the same two loops over
\(i\;\&\;(-i)\).  What changes is the \emph{interpretation} of slot \(i\): a
sum of \(A\), a sum of differences, one component of a paired encoding, a
count of ranks, a subtree's weight.  When a BIT problem feels hard, you are
almost always choosing a row of this table, not inventing a data structure.
\end{insightbox}

% ---------------- Taxonomy chart ----------------
\begin{figure}[H]
\centering
\begin{tikzpicture}[diagram, scale=0.95, transform shape,
    mnode/.style={circle, draw=#1, line width=0.9pt, fill=#1!8, minimum size=7mm,
                  font=\scriptsize\bfseries, text=#1!60!black}]
  \node[mnode=primaryblue] (a) at (0,0)     {lowbit};
  \node[annot, font=\scriptsize, anchor=west] at (1.25,0) {i\&(-i) = the block T[i] owns: [i-lowbit+1, i] (18.1)};
  \node[mnode=successgreen] (b) at (0,-1)  {walk};
  \node[annot, font=\scriptsize, anchor=west] at (1.25,-1) {query drops the last set bit; update adds it -- both <= log n steps (18.1)};
  \node[mnode=difforange]   (c) at (0,-2)  {difference};
  \node[annot, font=\scriptsize, anchor=west] at (1.25,-2) {store D[i]=A[i]-A[i-1]: range add becomes two point updates (18.2)};
  \node[mnode=warningred]   (d) at (0,-3)  {two BITs};
  \node[annot, font=\scriptsize, anchor=west] at (1.25,-3) {range-add/range-sum: prefix(x) = x*S1(x) - S2(x) over paired trees (18.2)};
  \node[mnode=softpurple]   (e) at (0,-4)  {lift};
  \node[annot, font=\scriptsize, anchor=west] at (1.25,-4) {find the first prefix >= k by walking powers of two down the tree (18.3)};
  \node[mnode=deepblue]     (f) at (0,-5)  {sweep};
  \node[annot, font=\scriptsize, anchor=west] at (1.25,-5) {sort queries on one axis, insert-and-count along the other; offline 2D and PBS (18.5)};
\end{tikzpicture}
\caption{The chapter's six-state taxonomy, read top to bottom as an
escalation.  Rows 1--2 are the machine; row~3 is a change of \emph{what you
store}; row~4 is a change of \emph{algebra} (two trees, one identity); row~5
runs the tree \emph{backwards} (search instead of sum); row~6 is a change of
\emph{schedule} -- queries arrive sorted, not online.  Every solved problem
sits on exactly one row: CSES~1648 on rows~1--2, CSES~1651 on row~3, the
two-BIT worked example on row~4, CF~61E and CSES~1143 on row~5, CF~383C on
row~3-plus-Ch.5, and the Summit CF~1181D spanning rows~5--6.}
\label{fig:c14-taxonomy}
\end{figure}

% ---------------- Who owns what ----------------
\subsection*{Who Owns What (so nothing is learned twice)}

\begin{center}
\small
\renewcommand{\arraystretch}{1.25}
\begin{tabular}{@{}p{0.30\linewidth} p{0.32\linewidth} p{0.32\linewidth}@{}}
\toprule
\textbf{Idea} & \textbf{Owner} & \textbf{This chapter's share} \\
\midrule
Euler tour, \(\tin/\tout\), ``subtree \(=\) interval'' & Ch.~5 (\chref{c24}) & cited by label; 18.4 supplies the range \emph{engine} Ch.~5 promised \\
segment trees, lazy propagation, arbitrary range ops & Ch.~19/20 (\chref{c13}, \chref{c15}) & boundary is prose only; BIT wins when the op is a sum/count and code size matters \\
merge-sort inversion counting & Ch.~1 (\chref{c01}) & same bound, offline divide\&conquer there; \emph{online} frequency sweep here (\secref{c14-order}) \\
LC 315 / LC 493 / LC 307, PBDS order-stat trees & Ch.~5 summit, Ch.~6 (\chref{c06}) & URLs live there; 18.3 is the array-index rival, mentioned in prose, never re-linked \\
stream median (two heaps) & Ch.~8 (\chref{c07}) & online medians stay there; 18.3 does k-th \emph{offline}/static \\
PBS scheduling, rollback DSU & Ch.~17 (\chref{c28}) & the lockstep idea is cited; 18.5 swaps the predicate to a BIT \\
\bottomrule
\end{tabular}
\end{center}

% ---------------- Recognition guide ----------------
\begin{recognitionbox}[Pattern card -- ``prefix structure under point changes, or counting across a sorted axis'']
\reckey{``update one position / sum a range'' arriving interleaved; ``add to a
whole interval, read one position''; ``how many elements \(\le k\) / the
k-th smallest'' with deletions; ``count pairs/triplets with index order and
value order''; ``for every subtree, add/read, and queries are many'';
``static 2D orthogonal counts, offline allowed''; ``the least value making a
sweep-count reach \(k\)'' (then lift the BIT inside a parallel binary
search).}
\recuse{
  \item point updates + prefix/range sums -- plain BIT, \secref{c14-core} (CSES~1648);
  \item whole intervals shift, single positions read -- difference BIT,
        \secref{c14-range} (CSES~1651); shifts \emph{and} interval sums -- two BITs;
  \item the answer is a \emph{rank} question (how many / k-th) over a
        compressed value range -- frequency BIT + lifting, \secref{c14-order} (CF~61E);
  \item the ranges are subtrees or root-paths of a tree -- Ch.5 flattening +
        any BIT flavour, \secref{c14-trees} (CF~383C);
  \item queries have a \emph{free} second dimension you can sort on, or a
        threshold you can binary-search offline -- \secref{c14-offline} (Summit).
}
\recnotuse{
  \item the range operation is not a sum/count (min with \emph{arbitrary}
        writes, gcd, max after decrements) -- that is a segment tree, Ch.~19
        (\chref{c13}); see \secref{c14-range}'s prefix-min caveat for the one
        order-preserving case;
  \item the tree ranges must be cut, linked or rerolled dynamically --
        Euler-tour \emph{trees}/splay live in Ch.~6 (\chref{c06});
  \item updates and queries interleave \emph{online} with no offline
        permission and you need full range-assign -- segment tree with lazy,
        Ch.~20 (\chref{c15});
  \item the whole array is static -- one prefix array is \(O(1)\) per query
        (Ch.~5, \propref{c24-prefix}); a BIT there is pure overhead.
}
One-line test: \emph{is the operation ``add at an index'' and the question
``sum/count a prefix'' -- possibly after sorting, possibly on a flattened
tree?  Then this chapter's array-index engine is the smallest structure that
answers it.}
\end{recognitionbox}

% ---------------- Constraint -> tool ----------------
\subsection*{Reading the Constraints (which flavour the limits want)}

\begin{center}
\small
\renewcommand{\arraystretch}{1.25}
\begin{tabular}{@{}p{0.30\linewidth} p{0.40\linewidth} p{0.24\linewidth}@{}}
\toprule
\textbf{Constraint signature} & \textbf{What it rules in/out} & \textbf{Tool} \\
\midrule
\(n,q\le2\cdot10^5\), point updates, range sums & prefix array dead (updates); BIT is the min-cost AC & \secref{c14-core} \\
range adds + point reads, \(u\) can be negative & point values drift far -- keep \code{long long} & \secref{c14-range} \\
values to \(10^9\), queries are rank/counting & no array of size \(10^9\): compress, then frequency BIT & \secref{c14-order} \\
\(n\le10^6\), ``count triplets'' with a middle index & two sweeps and a product; \(O(n\log n)\) with small constants only & \secref{c14-order} (CF 61E) \\
tree, subtree updates/reads, \(n,q\le2\cdot10^5\) & DFS order turns it into row 2 on slots \(1..n\) & \secref{c14-trees} \\
\(n,m,q\le5\cdot10^5\), ``which city gets the \(k\)-th shipment'' & simulate = TLE; sort + BIT + lifting answers each in \(\log n\) & \secref{c14-offline}, Summit \\
\bottomrule
\end{tabular}
\end{center}

% ---------------- Notation ----------------
\subsection*{Notation Used Only Here}

\(\operatorname{lowbit}(i)=i\,\&\,(-i)\) -- the value of the lowest set bit of
\(i\), i.e.\ the largest power of two dividing \(i\); \(T[i]\) is the BIT slot
covering the \term{responsibility interval} \([i{-}\operatorname{lowbit}(i){+}1,\,i]\);
\(\code{query}(i)=A_1+\cdots+A_i\) (a prefix sum, \(\mathit{pref}[i]\)) and
\(\code{update}(i,\delta)\) adds \(\delta\) to \(A_i\) everywhere it is
buried.  A subtree interval sum uses Ch.~5's reduction with this chapter's
\(1\)-based, \(\tout\) \textbf{inclusive} convention:
\(\sum_{x\in\subtree(v)}=\code{query}(\tout(v))-\code{query}(\tin(v)-1)\)
(Ch.~5's \(0\)-based half-open stamps compute the same number -- never mix
the two in one program; the pitfall box of \secref{c14-trees} shows the
symptom).  A \term{frequency BIT} indexes \emph{values} (after
\term{coordinate compression}: replace values by their rank in the sorted
set) and stores multiplicities; \term{lifting} denotes the search that walks
a candidate index down from the highest power of two, building the largest
prefix with sum \(<k\).  \term{Offline} means all queries are known upfront
and re-ordered before replay; \term{PBS} (parallel binary search, Ch.~17's
\secref{c28-offline}) runs one shared sweep per bit level for \emph{all}
queries at once.  Every section runs on \(A\) and \(T\) above: if a table
here and a trace there disagree, this opener is wrong and the section is
right -- fix it here, never the other way.
